\subsection{正元素}

设 \(G\) 是带预序关系 \(\preccurlyeq\) 的预序群；若 \(0 \preccurlyeq x\) 且 \(0 \preccurlyeq y\) ，则由(POM)推出 \(y \preccurlyeq x + y\) ，再由传递性得 \(0 \preccurlyeq x + y\) ；这表明集合 \(G_{+}\) ——其中的元素 \(x \in G\) 满足 \(0 \preccurlyeq x\) ——对加法封闭；此外，关系 \(x \preccurlyeq y\) 等价于 \(0 \preccurlyeq y - x\) ，即 \(y - x \in G_{+}\) 。反之：

命题3. --- 若P是阿贝尔群G的一个子集，它包含 0 且满足 \(P + P \subset P\) ，则关系 \(y - x \in P\) 是与G的群结构相容的预序关系。要使这个关系使G成为序群，必要且充分的是 \(P \cap (-P) = \{0\}\) ；要使G在此序下成为全序群，还须进而满足 \(P \cup (-P) = G\) 。

显然，关系 \(y - x \in P\) 是自反且传递的，并且（若把它写成 \(x \preccurlyeq y\) ）满足公理(POM)。为证明第二个断言，只需注意 \(P \cap (-P)\) 是子群 \(G'\) ，它由这样的元素 \(x\) 组成： \(0 \preccurlyeq x\) 且 \(x \preccurlyeq 0\) 。最后，说 \(G\) 是全序的意味着，对每一对元素 \(x, y\) （属于 \(G\) ）， \(x - y, y - x\) 这两个元素中至少有一个属于 \(P\) ，这就完成了证明。

定义3. --- 在序群G中，任何满足 \(0 \leqslant x\) （相应地， \(x \leqslant 0\) ）的元素 x 称为正元素（相应地，负元素）。

注意，0 是唯一既正又负的元素；任何满足 0 \textless{} x（相应地，x \textless{} 0）的元素称为严格正的（相应地，严格负的）。

例。 --- 在加法群 \(Z \times Z\) 中，设P是元素 \((x, y)\) 的集合，这些元素满足两个不等式 \(ax + by \geqslant 0\) 与 \(cx + dy \geqslant 0\) ，其中 a、b、c、d 是整数（*或实数*），且 \(ad - bc \neq 0\) ；“锥”P满足命题3的前两个条件。这样便可定义与 \(Z \times Z\) 的群结构相容的各种序；这个群在这些序下都不是全序的。

注记. --- 凭借条件 \(P + P \subset P\) ，在序群中关系 \(x \geqslant 0\) 蕴含对每个自然数 n 有 \(nx \geqslant 0\) 。此外，若G的正元素 x 具有有限阶 n，则 \(-x = (n - 1)x\) 是正的；因为

\[
\mathbf {P} \cap (- \mathbf {P}) = \{0 \},
\]

这蕴含 \(x = 0\) 。特别地，若 \(\mathbf{G}\) 的每个元素都具有有限阶，则 \(\mathbf{P} = \{0\}\) ；此时关系 \(x \leqslant y\) 等价于 \(x = y\) （离散序）。

